\clearpage
\SourceLeaf{019}\pdfbookmark[0]{Erstes Buch.}{liber-primus}
\begin{center}
{\Large ERSTES BUCH.}\par\bigskip
{\large ALLGEMEINE BEZIEHUNGEN ZWISCHEN DEN GRÖSSEN, DURCH DIE DIE BEWEGUNGEN DER HIMMELSKÖRPER UM DIE SONNE BESTIMMT WERDEN.}\par\bigskip
\pdfbookmark[1]{Erster Abschnitt.}{sectio-prima}
{\large ERSTER ABSCHNITT.}\par\medskip
\emph{Beziehungen für einen einzelnen Ort auf der Bahn.}\par\bigskip
\pdfbookmark[2]{1.}{article-001}1.
\end{center}
In diesem Werk werden wir die Bewegungen der Himmelskörper nur insoweit betrachten, als sie von der Anziehungskraft der Sonne bestimmt werden. Von unserer Untersuchung ausgeschlossen sind daher alle Nebenplaneten, die Störungen, welche die Hauptplaneten aufeinander ausüben, und jede Drehbewegung. Die bewegten Körper selbst betrachten wir als mathematische Punkte und setzen voraus, dass alle Bewegungen nach den folgenden Gesetzen vor sich gehen, die daher als Grundlage aller Untersuchungen in diesem Werk anzusehen sind.

I. Die Bewegung eines jeden Himmelskörpers erfolgt stets in ein und derselben Ebene, in der auch der Mittelpunkt der Sonne liegt.

II. Die vom Körper beschriebene Bahn ist ein Kegelschnitt mit einem Brennpunkt im Mittelpunkt der Sonne.

III. Die Bewegung auf dieser Bahn erfolgt so, dass die in verschiedenen
\SourceLeaf{020}Zeitintervallen um die Sonne überstrichenen Flächeninhalte diesen Intervallen proportional sind. Drückt man also die Zeiten und Flächeninhalte durch Zahlen aus, so ergibt jeder Flächeninhalt, geteilt durch die Zeit, in der die Fläche überstrichen wird, einen unveränderlichen Quotienten.

IV. Für verschiedene Körper, die sich um die Sonne bewegen, stehen die Quadrate dieser Quotienten im zusammengesetzten Verhältnis der den Bahnen entsprechenden Parameter und der jeweiligen Summen aus der Sonnenmasse und den Massen der bewegten Körper.

Bezeichnet man also mit $2p$ den Parameter der Bahn, auf der sich der Körper bewegt, mit $\mu$ die Materiemenge dieses Körpers (wobei die Masse der Sonne $=1$ gesetzt wird) und mit $\frac{1}{2}g$ den Flächeninhalt, den er in der Zeit $t$ um die Sonne überstreicht, so ist $\frac{g}{t\BareRoot{p}\cdot\BareRoot{(1+\mu)}}$ eine für alle Himmelskörper konstante Zahl. Da es gleichgültig ist, welchen Körper wir zur Bestimmung des Wertes dieser Zahl heranziehen, werden wir sie aus der Bewegung der Erde ableiten, deren mittleren Abstand von der Sonne wir als Längeneinheit wählen. Unsere Zeiteinheit ist stets der mittlere Sonnentag. Bezeichnet ferner $\pi$ das Verhältnis des Kreisumfangs zum Durchmesser, so ist der Flächeninhalt der ganzen von der Erde beschriebenen Ellipse offenbar $\pi\BareRoot{p}$. Er ist daher $=\frac{1}{2}g$ zu setzen, wenn man für $t$ das siderische Jahr einsetzt. Damit wird unsere Konstante~$=\frac{2\pi}{t\BareRoot{(1+\mu)}}$. Um den Zahlenwert dieser Konstanten zu ermitteln, die wir im Folgenden mit $k$ bezeichnen, setzen wir nach der neuesten Bestimmung das siderische Jahr, also $t=365{,}2563835$, und die Erdmasse, also $\mu=\frac{1}{354710}=0{,}000\,0028\,192$, an. Daraus ergibt sich

\begin{center}
\begin{tabular}{@{}ll@{}}
\makebox[15em][l]{$\log 2\pi$\hfill .\,.\,.\,.\,.\,.\,.\,.\,.\,.\,.\,.\,.\,.} & $0{,}798\,1798\,684$\\[.3em]
\makebox[15em][l]{Kompl. $\log t$\hfill .\,.\,.\,.\,.\,.\,.\,.\,.\,.} & $7{,}437\,4021\,852$\\[.3em]
\makebox[15em][l]{Kompl. $\log\BareRoot{(1+\mu)}$\hfill .\,.\,.\,.\,.} & $9{,}999\,9993\,878$\\[.3em]
\hline
\makebox[15em][l]{$\log k$\hfill .\,.\,.\,.\,.\,.\,.\,.\,.\,.\,.\,.\,.\,.\,.} & $8{,}235\,5814\,414$\\[.3em]
\multicolumn{1}{c}{$k=$} & $0{,}017\,2020\,9895$
\end{tabular}
\end{center}
